
\documentclass[12pt]{article}
\usepackage[utf8]{inputenc}
\usepackage[spanish]{babel}
\usepackage{amsmath}
\usepackage{amsfonts}
\usepackage{amssymb}
\usepackage{geometry}
\usepackage{listings}
\usepackage{xcolor}

\geometry{margin=2.5cm}

\title{Optimización - Universidad Nacional del Altiplano - FINESI\\
Evaluación: Unidad II}
\author{Apellidos y Nombres:Beatriz Umiña Machaca}
\date{}

\begin{document}

\maketitle

\section*{Pregunta 1}
\textbf{¿Cuáles son las principales diferencias entre la optimización de funciones convexas y no convexas, y cómo influyen en la elección de métodos?}\\


\subsection*{Respuesta:}
Las principales  diferencias son:

\textbf{Funciones Convexas:}
\begin{itemize}
    \item Cualquier mínimo local es también mínimo global
    \item Garantizan convergencia a la solución óptima
    \item Métodos eficientes: descenso de gradiente, Newton, quasi-Newton
    \item Condiciones de optimalidad: $\nabla f(x^*) = 0$ es condición necesaria y suficiente
\end{itemize}

\textbf{Funciones No Convexas:}
\begin{itemize}
    \item Múltiples mínimos locales
    \item No garantizan convergencia al mínimo global
    \item Requieren métodos más sofisticados: algoritmos genéticos, simulated annealing, métodos de enjambre
    \item Condiciones de optimalidad: $\nabla f(x^*) = 0$ solo es condición necesaria
\end{itemize}

\textbf{Influencia en la elección de métodos:}
Para funciones convexas se prefieren métodos determinísticos y eficientes, mientras que para no convexas se requieren métodos estocásticos o metaheurísticos que exploren mejor el espacio de búsqueda.

\section*{Pregunta 2}
\textbf{¿Cuándo resulta más efectivo usar variantes del descenso de gradiente y cuáles son los factores clave en su selección?}

\subsection*{Respuesta:}

\textbf{Variantes y sus casos de uso:}
\begin{itemize}
    \item \textbf{Gradient Descent (GD):} Datasets pequeños, funciones suaves
    \item \textbf{Stochastic Gradient Descent (SGD):} Datasets grandes, cuando se requiere velocidad
    \item \textbf{Mini-batch GD:} Balance entre precisión y eficiencia computacional
    \item \textbf{Adam:} Redes neuronales, problemas con gradientes dispersos
    \item \textbf{RMSprop:} Redes neuronales recurrentes
    \item \textbf{Adagrad:} Problemas con características dispersas
\end{itemize}

\textbf{Factores clave en la selección:}
\begin{enumerate}
    \item Tamaño del dataset
    \item Naturaleza del problema (convexo/no convexo)
    \item Recursos computacionales disponibles
    \item Velocidad de convergencia requerida
    \item Estabilidad del gradiente
    \item Características del espacio de parámetros
\end{enumerate}

\section*{Pregunta 3}
\textbf{¿Cómo se aplican los criterios de optimalidad en un método de optimización?}

\subsection*{Respuesta:}

\textbf{Criterios de Optimalidad:}

\textbf{Condiciones de primer orden (necesarias):}
$$\nabla f(x^*) = 0$$

\textbf{Condiciones de segundo orden (suficientes):}
$$\nabla^2 f(x^*) \succ 0 \text{ (definida positiva)}$$

\textbf{Aplicación en métodos de optimización:}
\begin{enumerate}
    \item \textbf{Criterio de parada:} $\|\nabla f(x_k)\| < \epsilon$
    \item \textbf{Dirección de búsqueda:} $d_k = -\nabla f(x_k)$
    \item \textbf{Verificación de convergencia:} Evaluar si se cumple $\nabla f(x_k) \approx 0$
    \item \textbf{Test de optimalidad:} Verificar que la matriz Hessiana sea definida positiva
\end{enumerate}

\textbf{Con restricciones (Condiciones KKT):}
\begin{align}
\nabla f(x^*) + \sum_{i=1}^m \lambda_i \nabla g_i(x^*) &= 0\\
g_i(x^*) &\leq 0, \quad i = 1,\ldots,m\\
\lambda_i &\geq 0, \quad i = 1,\ldots,m\\
\lambda_i g_i(x^*) &= 0, \quad i = 1,\ldots,m
\end{align}

\section*{Pregunta 4}
\textbf{¿De qué manera facilitan las herramientas de optimización automática (p.ej., Optuna) la búsqueda de configuraciones óptimas en aprendizaje automático?}

\subsection*{Respuesta:}

\textbf{Facilidades que proporciona Optuna:}

\begin{itemize}
    \item \textbf{Algoritmos de optimización avanzados:} TPE (Tree-structured Parzen Estimator), CMA-ES, búsqueda aleatoria
    \item \textbf{Pruning automático:} Detiene trials no prometedores tempranamente
    \item \textbf{Paralelización:} Ejecución distribuida de múltiples trials
    \item \textbf{Persistencia:} Almacenamiento de estudios y reanudación
    \item \textbf{Visualización:} Gráficos de convergencia y análisis de parámetros
\end{itemize}

\textbf{Ventajas en ML:}
\begin{enumerate}
    \item Automatización completa del proceso de búsqueda
    \item Manejo eficiente de espacios de hiperparámetros complejos
    \item Reducción significativa del tiempo de búsqueda
    \item Integración sencilla con frameworks de ML
    \item Soporte para optimización multiobjetivo
\end{enumerate}

\section*{Pregunta 5}
\textbf{¿Cómo influye la selección de estrategias de regularización en el desempeño y la generalización de un modelo?}

\subsection*{Respuesta:}

\textbf{Tipos de regularización y sus efectos:}

\textbf{L1 (Lasso):}
$$J(w) = \frac{1}{2m}\sum_{i=1}^m (h_w(x^{(i)}) - y^{(i)})^2 + \lambda\sum_{j=1}^n |w_j|$$
\begin{itemize}
    \item Produce sparsity (selección de características)
    \item Útil cuando hay muchas características irrelevantes
\end{itemize}

\textbf{L2 (Ridge):}
$$J(w) = \frac{1}{2m}\sum_{i=1}^m (h_w(x^{(i)}) - y^{(i)})^2 + \lambda\sum_{j=1}^n w_j^2$$
\begin{itemize}
    \item Reduce la magnitud de los pesos
    \item Mejor para multicolinealidad
\end{itemize}

\textbf{Elastic Net:}
$$J(w) = \frac{1}{2m}\sum_{i=1}^m (h_w(x^{(i)}) - y^{(i)})^2 + \lambda_1\sum_{j=1}^n |w_j| + \lambda_2\sum_{j=1}^n w_j^2$$

\textbf{Dropout:}
\begin{itemize}
    \item Desactiva aleatoriamente neuronas durante el entrenamiento
    \item Previene co-adaptación de características
\end{itemize}

\textbf{Impacto en el desempeño:}
\begin{itemize}
    \item Reduce overfitting
    \item Mejora la generalización
    \item Puede reducir la capacidad del modelo si se aplica excesivamente
    \item Requiere ajuste cuidadoso del parámetro de regularización $\lambda$
\end{itemize}

\section*{Pregunta 6}
\textbf{¿Cómo se pueden integrar rutinas de validación y corrección de código en Optuna para mejorar la calidad y eficacia del proceso de optimización de hiperparámetros?}

\subsection*{Respuesta:}

\textbf{Estrategias de integración:}

\begin{enumerate}
    \item \textbf{Validación cruzada dentro del objetivo:}
    \begin{itemize}
        \item Usar k-fold cross-validation en cada trial
        \item Reportar métricas estadísticamente robustas
    \end{itemize}
    
    \item \textbf{Manejo de excepciones:}
    \begin{itemize}
        \item Try-catch blocks para capturar errores
        \item Retornar valores predeterminados en caso de fallo
    \end{itemize}
    
    \item \textbf{Callbacks personalizados:}
    \begin{itemize}
        \item Validación de parámetros antes de entrenar
        \item Logging detallado de errores
    \end{itemize}
    
    \item \textbf{Pruning inteligente:}
    \begin{itemize}
        \item Usar MedianPruner o PercentilePruner
        \item Configurar early stopping basado en métricas de validación
    \end{itemize}
\end{enumerate}

\textbf{Ejemplo de implementación:}
\begin{lstlisting}[language=Python]
def objective(trial):
    try:
        # Validación de parámetros
        params = {
            'learning_rate': trial.suggest_float('lr', 1e-5, 1e-1, log=True),
            'batch_size': trial.suggest_categorical('batch_size', [16, 32, 64])
        }
        
        # Validación cruzada
        cv_scores = cross_val_score(model, X, y, cv=5)
        
        # Reporte de métricas intermedias
        for i, score in enumerate(cv_scores):
            trial.report(score, i)
            if trial.should_prune():
                raise optuna.TrialPruned()
                
        return cv_scores.mean()
    
    except Exception as e:
        logging.error(f"Error in trial {trial.number}: {e}")
        return float('inf')  # Penalizar trials con errores
\end{lstlisting}

\section*{Pregunta 7}
\textbf{¿Cómo podemos utilizar las métricas provistas por scikit-learn para evaluar un modelo de clasificación, y qué información adicional brindan además de la exactitud?}

\subsection*{Respuesta:}

\textbf{Código completado:}
\begin{lstlisting}[language=Python]
from sklearn.metrics import accuracy_score, precision_score, recall_score, f1_score, classification_report, confusion_matrix

y_true = [0, 1, 2, 2, 2]
y_pred = [0, 2, 2, 2, 1]

print("Accuracy:", accuracy_score(y_true, y_pred))
print("Precision:", precision_score(y_true, y_pred, average='macro'))
print("Recall:", recall_score(y_true, y_pred, average='macro'))
print("F1-score:", f1_score(y_true, y_pred, average='macro'))
print("Classification Report:")
print(classification_report(y_true, y_pred))
print("Confusion Matrix:")
print(confusion_matrix(y_true, y_pred))
\end{lstlisting}

\textbf{Información adicional de las métricas:}

\begin{itemize}
    \item \textbf{Precision:} $\frac{TP}{TP + FP}$ - Proporción de predicciones positivas correctas
    \item \textbf{Recall (Sensibilidad):} $\frac{TP}{TP + FN}$ - Proporción de casos positivos correctamente identificados
    \item \textbf{F1-score:} $2 \cdot \frac{Precision \cdot Recall}{Precision + Recall}$ - Media armónica entre precisión y recall
    \item \textbf{Specificity:} $\frac{TN}{TN + FP}$ - Proporción de casos negativos correctamente identificados
\end{itemize}

\textbf{Ventajas sobre accuracy:}
\begin{enumerate}
    \item Mejor evaluación en datasets desbalanceados
    \item Información específica por clase
    \item Identificación de sesgos del modelo
    \item Comprensión del tipo de errores cometidos
\end{enumerate}

\textbf{Interpretación de resultados:}
Con los datos dados:
\begin{itemize}
    \item Accuracy = 0.6 (60\% de predicciones correctas)
    \item La precision y recall por clase revelan el desempeño específico
    \item La matriz de confusión muestra exactamente qué clases se confunden
\end{itemize}

\end{document}
```